\documentclass[11pt,a4paper]{article}
\usepackage[utf8]{inputenc}
\usepackage[T1]{fontenc}
\usepackage[spanish]{babel}
\usepackage{amsmath,amssymb,mathtools}
\usepackage{geometry}
\usepackage{enumitem}
\usepackage{lmodern}
\usepackage{microtype}
\geometry{margin=2.4cm}
\setlength{\parindent}{0pt}
\setlength{\parskip}{0.5em}
\newcommand{\R}{\mathbb{R}}
\newcommand{\Span}{\operatorname{L}}
\newcommand{\Ima}{\operatorname{Im}}
\newcommand{\sen}{\operatorname{sen}}
\begin{document}

\begin{center}
{\LARGE\bfseries Álgebra Lineal}\\[0.4em]
{\Large Tema 4: El Espacio Afín --- Ejercicios}
\end{center}

\section{Ejercicio 1}
Dados $R=\{O,u_1,u_2\}$ y $R'=\{O',v_1,v_2\}$, encontrar las fórmulas del cambio de sistema de referencia sabiendo que:
\[
[\overrightarrow{OO'}]_B=(2,3),
\qquad
[v_1]_B=(1,-1),
\qquad
[v_2]_B=(5,4).
\]

\section{Ejercicio 2}
Demostrar que el conjunto de soluciones de un sistema lineal compatible:
\[
AX=C
\]
es un espacio afín de $\R^n$.

\section{Ejercicio 3}
En $\R^4$:
\[
F_1=
\left\{
\begin{array}{l}
x_1=a+3\lambda+2\mu\\
x_2=1-\lambda-\mu\\
x_3=4+\lambda\\
x_4=6+5\lambda+2\mu
\end{array}
:
\lambda,\mu\in\R
\right\},
\]
\[
F_2=
\left\{
\begin{array}{l}
x_1=2+\alpha+2\beta\\
x_2=1\\
x_3=1+\alpha+\beta\\
x_4=3\alpha
\end{array}
:
\alpha,\beta\in\R
\right\}.
\]
\begin{enumerate}[label=\alph*)]
\item Calcular $a$ para que $F_1\cap F_2\neq\emptyset$.
\item Para dicho valor, obtener $F_1\cap F_2$ y $F_1+F_2$.
\end{enumerate}

\section{Ejercicio 4}
Respecto de $R=\{O,e_1,e_2,e_3\}$:
\[
A=(3,1,-2),\quad
B=(2,2,0),\quad
C=(1,0,-1),\quad
D=(4,3,-2).
\]
Las coordenadas de $P$ en:
\[
R'=\{A;\overrightarrow{AB},\overrightarrow{AC},\overrightarrow{AD}\}
\]
son $(x',y',z')$. Obtener las coordenadas de $P$ en $R$.

\section{Ejercicio 5}
Dadas:
\[
x_1=1+y_2,\qquad
x_2=\frac13y_1-\frac23y_3,\qquad
x_3=2+\frac23y_1+\frac13y_3,
\]
donde $x_i$ son coordenadas en $R$ e $y_i$ en $R'$:
\begin{enumerate}[label=\alph*)]
\item Obtener la expresión matricial del cambio.
\item Obtener la ecuación de $x_1+x_2+x_3=3$ en $R'$.
\item Dar las ecuaciones paramétricas e implícitas, punto y espacio de dirección en $R$ y $R'$.
\end{enumerate}

\section{Ejercicio 6}
En $\R^4$:
\[
F_1=\{(x_1,x_2,x_3,x_4):x_1+x_2=4,\;x_3+x_4=a\},
\]
\[
F_2=\{(3+\alpha,2-2\alpha,2\alpha,-1+\alpha):\alpha\in\R\}.
\]
Encontrar $a$ para que la variedad afín generada por $F_1$ y $F_2$ tenga dimensión mínima y calcularla.

\section{Ejercicio 7}
Sea:
\[
V=\{p(x)\in\mathcal P_2(x):p(1)=1,\;p(2)=2\}.
\]
\begin{enumerate}[label=\alph*)]
\item Demostrar que $V$ es una variedad afín.
\item Obtener un sistema de referencia afín de $V$ y las coordenadas de $2x^2-5x+4$.
\item Respecto de $R=\{0;1,x,x^2\}$, obtener el hiperplano que pasa por $x$ y es perpendicular a $V$.
\end{enumerate}

\section{Ejercicio 8}
Determinar la aplicación afín $f:\mathcal A\to\mathcal A'$ tal que:
\[
f(1,2)=(1,2,3),
\]
\[
\widetilde f(e_1)=e_1+4e_2,
\qquad
\widetilde f(e_2)=e_1-e_2+e_3.
\]
\begin{enumerate}[label=\alph*)]
\item Escribir matricialmente las ecuaciones de $f$.
\item Obtener las ecuaciones paramétricas de $\Ima f$ y su espacio de dirección.
\end{enumerate}

\section{Ejercicio 9}
Dada:
\[
C\equiv(x_1-1)^2+(x_2-1)^2=9,
\]
obtener su ecuación en:
\[
R'=\{A,u_1,u_2\},
\quad
A=(3,3),
\]
\[
u_1=\left(\frac1{\sqrt2},\frac1{\sqrt2}\right),
\qquad
u_2=\left(\frac2{\sqrt2},-\frac2{\sqrt2}\right).
\]
¿Qué tipo de curva es $C$ en $R'$?

\section{Ejercicio 10}
En $\R^3$:
\[
\pi:x_1+x_2+x_3=2.
\]
Sea:
\[
R'=\{P,u_1,u_2,u_3\},
\quad
P=(1,1,0),
\]
\[
u_1=-e_1+e_3,\qquad
u_2=-e_2+e_3,\qquad
u_3=e_1+e_2+e_3.
\]
Calcular la ecuación de $\pi$ respecto de $R'$.

\section{Ejercicio 11}
En $\R^3$:
\[
A_1=(1,0,0)+\Span\{e_1+e_2+e_3\},
\]
\[
A_2=(0,0,0)+\Span\{2e_1+e_2+2e_3\}.
\]
Calcular $\dim(A_1+A_2)$ y relacionarla con $\dim A_1$ y $\dim A_2$.

Si:
\[
A_3=\{(x,y,z):2x-y-z=0\},
\]
calcular $\dim(A_1+A_3)$ y relacionarla con $\dim A_1$ y $\dim A_3$.

\section{Ejercicio 12}
Determinar cuáles son aplicaciones afines:
\[
\text{a) }f(x,y)=(x+1,y-1)
\]
\[
\text{b) }f(x,y)=(2x,2y)
\]
\[
\text{c) }f(x,y)=(2x-1,3y+2)
\]
\[
\text{d) }f(x,y)=(3x+1,0)
\]
\[
\text{e) }f(x,y)=(x^2+1,y)
\]

\section{Ejercicio 13}
Encontrar la expresión analítica del giro de centro $(1,1)$ y ángulo $\pi/2$.

\section{Ejercicio 14}
Encontrar la expresión analítica de la proyección ortogonal sobre:
\[
x+y=1.
\]

\section{Ejercicio 15}
Encontrar la expresión analítica de la simetría ortogonal respecto de:
\[
x+y=1.
\]

\section{Ejercicio 16}
En $\R^3$:
\[
F=\{(x_1,x_2,x_3):x_1+x_2-x_3=1\}.
\]
\begin{enumerate}[label=\alph*)]
\item Analizar que $F$ es una variedad afín e indicar sus elementos.
\item Si:
\[
R'=\{P,u_1,u_2,u_3\},
\quad
P=(-2,1,3),
\]
\[
u_1=e_1-e_3,\qquad
u_2=e_2-e_3,\qquad
u_3=e_1+e_2+e_3,
\]
determinar la ecuación de $F$ respecto de $R'$.
\end{enumerate}

\end{document}
